\section{Vectores y matrices.}
Objetivos:
\begin{itemize}
    \item Familiarizarse con los conceptos de vectores y matrices.
    \item Realizar operaciones de suma y multiplicación por escalar.
    \item Estudiar la definición multiplicación entre matrices y vectores.
    \item Realizar operaciones de la multiplicación entre matrices y
vectores.
\end{itemize}

\subsection{Definiciones generales}
Un \textbf{vector fila} de $n$ componentes (entradas) se define como un conjunto ordenado de $n$ números dispuestos de manera horizontal, es decir 
\[x= \left(\begin{array}{ccccc}
   x_1 & x_2 &  \dots&  x_n 
\end{array}\right),\]
donde cada componente $x_i$ pertenece a un conjunto $K$, usualmente $\R$ o $\C$. 
De la misma manera, un \textbf{vector columna} de $n$ componentes es un conjunto ordenado de $n$ números escritos de forma vertical
\[ x=\left(\begin{array}{c}
   x_1 \\
   x_2 \\
   \vdots\\
   x_n 
\end{array}\right).\]
En ambos casos, si un vector tiene $n$ componentes, se dice que el vector es de dimensión $n$.
\begin{defin}
Una \textbf{matriz $A$ de $m \times n$} es un arreglo rectangular de $mn$ elementos de un conjunto dispuestos en $m$ filas y $n$ columnas:
\[A=\left(\begin{array}{cccccc}
     a_{11}& a_{12} & \dots & a_{1j} &\dots & a_{1n}\\
     a_{21}& a_{22} & \dots &a_{2j} &\dots & a_{2n}\\
     \vdots&  \vdots & \ddots &  \vdots &\ddots &  \vdots\\
     a_{i1}& a_{i2} & \dots &a_{ij} &\dots & a_{in}\\
     \vdots&  \vdots & \ddots & \vdots &\ddots &  \vdots\\
     a_{m1}& a_{m2} & \dots & a_{mj} &\dots & a_{mn}
\end{array}\right)_{m\times n}.\]





De forma abreviada, la matriz $A$ puede denotarse por $A=(a_{ij})$, con $i=1,2,\dots, m$, $j=1,2,\dots, n$, donde el elemento $a_{ij}$ corresponde al número que aparece en la fila $i$ y la columna $j$ de la matriz $A$.

Las matrices usualmente se denotan por letras mayúsculas $A, B,$ etc. y sus respectivas entradas se denotan por la minúscula respectiva $a_{ij}, b_{ij}$, etc.

Si $m =n$ se dice \textbf{matriz cuadrada} y en caso contrario se dice que la matriz es \textbf{rect\'angular}.

El conjunto de matrices de $m\times n$ con entradas en el cuerpo $K$ se denota por $\M_{m\times n}(K)$, en otras palabras
\[\M_{m\times n}(K):=\Big\{A=(a_{ij}) \colon a_{ij}\in K \text{ con $i=1,2,\dots, m$ y $j=1,2,\dots, n$}\Big\}.\]
\end{defin}
En el caso de matrices cuadradas de orden $n$ es escribe $\M_n(K)$.


Una matriz $A$ de $m \times n$ cuyas entradas son todas  iguales a 0, se denomina \textbf{matriz cero} de $m \times n$ y se denota por $0_{m\times n}$.

La matriz cuadrada $I_n=(a_{ij})$ de $n\times n$ tal que 
\[a_{ij}=\left\{\begin{array}{ll}
    1 & \text{ si } i= j, \\
   0 & \text{ si no},
\end{array}\right.\]
se llama \textbf{matriz identidad} de orden $n$. Su nombre se debe a sus propiedades aritm\'eticas, que discutiremos m\'as adelante.

Dos matrices $A=(a_{ij})$ y $B=(b_{ij})$ son iguales si son del mismo tamaño y las componentes correspondientes son iguales, es decir, $a_{ij}=b_{ij}$ para cada $i$ y para cada $j$.

Si $A = (a_{i j})$ es una matriz de orden $m \times n$, la \textbf{transpuesta} de $A$, denotada por $A^t$,
es la matriz de orden $n \times m$ que se obtiene al intercambiar las filas por las columnas
de $A$. En consecuencia, 
\[A^t := (a_{j i }).\]

\begin{ejem}
\begin{enumerate}[(a)]
    \item Un vector fila de dimensi\'on $n$ es una matriz de $1\times n$.
    \item Un vector columna de dimensi\'on $n$ es una matriz de $n\times 1$.
    \item $A=\left(\begin{array}{cc}{} 1 & 2\\ 3 & 4 \end{array}\right)$ es una matriz cuadrada de $2\times 2$.
    \item Las matrices $\begin{pmatrix}1& 6& 2 \\0&3&4\end{pmatrix}_{2\times 3}$ y $\left(\begin{array}{cc}
   1& 0 \\
    6&3\\
    2& 4
\end{array}\right)_{3\times 2}$ son distintas.
    \item La matriz $\begin{pmatrix}
   0& 0& 0& 0 \\
   0&0&0& 0
\end{pmatrix}_{2\times 4}$ es la matriz cero de orden $2\times 4$.
\end{enumerate}
\end{ejem}

\begin{ejer}
\begin{enumerate}[(a)]
    \item Determine la matriz de $4\times 3$ tal que $a_{ij}=i-j$.
    \item De un ejemplo de una matriz $A=(a_{ij})$ tal que $a_{ij}=a_{ji}$ para todo $i$ y para todo $j$.
    \item ?`Existe alguna relaci\'on entre $\M_{1\times n}(\R)$ y $\R^n$?
    \item Construya la matriz $A = (a_{ij})$ de $3 \times 4$, donde
\[a_{ij}=\left\{\begin{array}{ll}
    i & \text{ si } i\leq j, \\
    i+j & \text{ si } i>j
\end{array}\right.\]
\item Encuentre los n\'umeros reales $x,y,z,u,v$ y $w$ tales que se cumple la igualdad
\[\begin{pmatrix}
   -1& x& x-y \\
  u&3&-1
\end{pmatrix}=\begin{pmatrix}
   z& 4& 7 \\
  2&x+v&u+w
\end{pmatrix}\]
Respuesta: $x = 4, y = -3, z = -1, u = 2, v = -1, w = -3$.
\end{enumerate}
\end{ejer}

\begin{rmk}\textbf{?`Por qu\'e estudiar matrices?}

Por su simplicidad y sus propiedades algebraicas, las matrices pueden utilizarse en diversas situaciones, como ordenar datos o representar situaciones pr\'acticas.

La producción semanal, en cientos, de los artículos $p_1$, $p_2$,\dots ,$p_{60}$ que se fabrican en una industria se pueden expresar mediante una matriz $P$ de 60 filas - donde se escribirán los productos elaborados - y 5 columnas que indicarán los días de la semana de lunes a viernes:
\[
P=\begin{blockarray}{cccccc}
    Lu & Ma & Mi &Ju & Vi \\
\begin{block}{(ccccc)c}
 2& 3,5 & 3 &3,2  & 3 & p_1 \\
  6 & 6,5 & 6,3 & 6,2 & 6 & p_2 \\
  \vdots & \vdots & \vdots & \vdots & \vdots & \vdots \\
  4 & 3,8 & 3,5 & 4 & 3,2 & p_{60} \\
\end{block}
\end{blockarray}
 \]
\end{rmk}

\begin{defin}
La \textbf{diagonal de una matriz} $A$ son las entradas $a_{i j}$, tales que $i=j$ y se denota por $\diag(A):=(a_{11}, \dots, a_{nn})$.

Para una matriz cuadrada de orden $n$, se define la \textbf{traza de $A$} como
\[\tr(A):=\sum_{i=1}^{n} a_{ii},\]
es decir, la traza corresponde a la suma de los elementos de la diagonal de una matriz.

Mientras que para un vector $(a_{1} ,a_{2} ,  \dots  a_{n})$ se define \textbf{la matriz diagonal} a la matriz cuadrada
\[\diag(a_{1} ,a_{2} ,  \dots  a_{n}):=\left\{\begin{array}{ll}
    a_i & \text{ si } i=j,\\
    0 & \text{ si no.}
\end{array}\right.\]
\end{defin}


\begin{ejem}
\begin{enumerate}
    \item     $A=\begin{pmatrix}
    2 & 0 & 0\\
    0 & -5 & 0\\
    0 & 0 & 1\\
    \end{pmatrix}
    $ y $B=\begin{pmatrix}
    -1 & 0 \\
    0 & 0 \\
    \end{pmatrix}$ son matrices diagonales.
    \item Para $n\in \N$ se tiene que $\tr(I_n)=n$.
\end{enumerate}

\end{ejem}


\begin{defin}
Una matriz cuadrada $A = (a_{i j})$ se llama \textbf{triangular
superior} si $a_{i j} = 0$ para $i > j$, y se llama \textbf{triangular inferior}
si $a_{i j} = 0$ para $i < j$.

\end{defin}

\begin{ejem}
    $C=\begin{pmatrix}
    4 & -1 & 5\\
    0 & -6 & 0\\
    0 & 0 & 8\\
    \end{pmatrix}
    $ y $D=\begin{pmatrix}
    -3 & 0 \\
    2 & 7 \\
    \end{pmatrix}$ son matrices triangulares superior e inferior respectivamente.
\end{ejem}

\begin{ejem}
    \begin{enumerate}[(a)]
        \item Determine por extensi\'on la matriz $A=(a_{ij})$ de orden 3 dada por
        \[a_{ij}=\begin{cases}
			i, & \text{si $i=j$,}\\
            2i-j, & \text{si $i<j$},\\
            |i-3j| & \text{si $i>j$}.
		 \end{cases}\]
		 \item ¿Cuál es la diagonal y la traza de $A$?
        \item ¿Si la matriz $A$ fuese de orden 20? ¿y si fuese de orden $n$?
    \end{enumerate}
\end{ejem}
% \begin{sol}
% \begin{enumerate}[(a)]
%     \item Directo de la definici\'on se tiene que la matriz 
%     \[A=\begin{pmatrix}
%     a_{11} & a_{12} & a_{13}\\
%     a_{21} & a_{22} & a_{23}\\
%     a_{31}& a_{32} & a_{33}\\
%     \end{pmatrix}=\begin{pmatrix}
%     1 & 0 & -1\\
%     1 & 2 & 1\\
%     0& 3 & 3\\
%     \end{pmatrix},\]
% luego su diagonal es $\diag(A)=(1,2,3)$ mientras que $\tr(A)=6$.  
% \item  Si $A$ es de orden 20 entonces $\tr(A) =\sum_{i=1}^{20} a_{ii}=\sum_{i=1}^{20} i=210$

% \item Si la matriz es de orden $n$, tenemos que $\tr(A) =\sum_{i=1}^{n} a_{ii}=\sum_{i=1}^{n} i=\dfrac{n(n+1)}{2}.$
% \end{enumerate}
% \end{sol}
\begin{defin}
Sea $A = (a_{i j})$ matriz cuadrada. Se dice que $A$ es una matriz \textbf{simétrica} si
$A^t = A$ y se dice que $A$ es \textbf{antisimétrica} si $A^t = -A$, donde $-A$ es la matriz
$-A = (-a_{i j})$.
\end{defin}

\begin{ejem}
    La matriz $A=\begin{pmatrix}
    6 & -1 & 4\\
    -1 & 7 & 2\\
    4& 2 & -3\\
    \end{pmatrix}$ es una matriz simétrica.
\end{ejem}
\begin{ejer}
\begin{enumerate}[(a)]
\item Construya una matriz antisimétrica de orden 4.
\item Si $A$ es una matriz de orden $n$ antisimétrica, demuestre que $ \diag(A)=(0,...,0)$ y que
$\tr(A)=0$. 
\end{enumerate}
\end{ejer}

\begin{ejer}
¿Cuál es la transpuesta de la matriz $A = (a_{ij})$ de orden $3 \times 2$ definida por $a_{i j} = 3i-j^2?$
\end{ejer}

\begin{prop}\label{traspuesta}
$(A^t)^t=A$ para toda $A\in M_{m\times n}(K)$.
\end{prop}
\proof (Ejercicio)



\subsection{Adici\'on y multiplicaci\'on de matrices por un escalar.}
En el conjunto de matrices $M_{m\times n}(K)$ es posible definir dos operaciones:
\begin{defin}
Sean $A = (a_{ij})$ y $B = (b_{ij})$ dos matrices de $m \times n$. Entonces la suma
de $A$ y $B$ es la matriz de $m \times n$
\[A+B=(c_{ij}):=\left(\begin{array}{ccccc}
     a_{11}+b_{11} & \dots & a_{1j}+b_{1j} &\dots & a_{1n}+b_{1n}\\
     a_{21}+b_{21} & \dots &a_{2j}+b_{2j} &\dots & a_{2n}+b_{2n}\\
     \vdots&  \vdots & \ddots &  \vdots &\vdots \\
     a_{m1}+b_{m1} & \dots & a_{mj}+b_{mj} &\dots & a_{mn}+b_{mn}
\end{array}\right)_{m\times n},\]
es decir, $c_{ij}=a_{ij}+b_{ij}$ para cada $i,j$.
\end{defin}

\begin{ejem}
\begin{align*}
    \begin{pmatrix} 1 & 4 & 0 \\ -2 & 6 & 5 \end{pmatrix} + \begin{pmatrix} -3 & 1 & -1 \\ 3 & 0 & 2 \end{pmatrix} &= \begin{pmatrix} 1-3 & 4+1 & 0 -1 \\ -2+3 & 6 + 0 & 5+2 \end{pmatrix}\\
&= \begin{pmatrix} -2 & 5 & -1 \\ 1 & 6 & 7 \end{pmatrix}
\end{align*}
\end{ejem}

\begin{rmk}
La suma $\begin{pmatrix} 1 & 4 & 0 \\ -2 & 6 & 5 \end{pmatrix} + \begin{pmatrix} 4 & 3 \\ 2 & 1 \end{pmatrix} $ 
\\\
no está definida pues las matrices son de distinto tamaño.
\end{rmk}

\begin{defin}
Sea $G$ un conjunto con una operaci\'on (binaria)
\begin{align*}
* \colon G\times G &\to G \\
(a,b) &\mapsto a*b
\end{align*}
diremos que $(G,*)$ es un grupo conmutativo(abeliano), si satisface los siguientes \textit{axiomas}:
\begin{enumerate}
    \item Asociatividad: $(a*b)*c=a*(b*c)$ para todo $a,b,c\in G$.
    \item Elemento neutro: Existe un (\'unico) elemento $e\in G$ tal que para todo $a\in G$ se tiene que $e*a=a*e=a$. Dicho elemento se llama \textbf{elemento neutro del grupo $G$}.
    \item Inverso: Para cada $a\in G$ existe un (\'unico) elemento $b\in G$ tal que $a*b=b*a=e$, donde $e$ es la identidad del grupo. Para cada $a$ el elemento que cumpla la propiedad se llama \textbf{inverso de $a$}.
    \item Conmutatividad: $a*b=b*a$ para todo $a,b\in G$.
\end{enumerate}
\end{defin}

\begin{thm}
El conjunto $\M_{m\times n}(K)$ dotado de la operaci\'on suma, es un grupo abeliano, cuyo elemento neutro es $0_{m\times n}$ y el elemento inverso de una matriz $A=(a_{ij})$ es $-A=(-a_{ij})$
\end{thm}

Se define $A - B := A + (-B)$, para todo $ A, B \in \M_{m\times n}(K)$.

\begin{thm}
Se tiene las siguientes propiedades:
\begin{enumerate}
    \item $\tr(A+B) = \tr(A) + \tr(B)$, para todo $ A, B \in \M_{m\times n}(K)$.
    \item $(A + B)^t = A^t + B^t$, para todo $ A, B \in \M_{m\times n}(K)$.
    \item $A, B$ diagonales, entonces $A + B$ es diagonal.
    \item $A, B$ simétricas, entonces $A + B$ es simétrica.
    \item $A, B$ antisimétricas, entonces $A + B$ es antisimétrica.
    \item $- ( -A) = A$, para toda $ A \in \M_{m\times n}(K)$.
    \item $-(A + B) = -A + (-B)$, para todo $ A, B \in \M_{m\times n}(K)$.
    \item $ A + C = B + C$, entonces  $A = B$.
    \item $ A + X = B$, entonces $X = B - A$.
\end{enumerate}
\end{thm}
\proof (Ejercicio)


\begin{defin}
Si $A$ es una matriz de $m \times n$ y si $\alpha$ es un escalar, es decir un elemento cualquiera del conjunto $K$, entonces la matriz $\alpha \cdot A$ está dada por
\[\alpha \cdot A:=(\alpha \cdot a_{ij})\left(\begin{array}{ccccc}
     \alpha a_{11} & \dots &\alpha  a_{1j} &\dots &\alpha  a_{1n}\\
     \alpha a_{21} & \dots &\alpha  a_{2j} &\dots &\alpha  a_{2n}\\
     \vdots&  \vdots & \ddots &  \vdots &\vdots \\
     \alpha a_{m1} & \dots &\alpha  a_{mj} &\dots &\alpha  a_{mn}
\end{array}\right)_{m\times n}.\]
En general omitiremos el punto ($\cdot$), escribiendo $\alpha A$ en vez de $\alpha \cdot A$.
\end{defin}

\begin{ejem}
    $3\begin{pmatrix}
    2 & -1 & 4\\
    -1 & 5 & 2\\
    \end{pmatrix}=\begin{pmatrix}
    6 & -3 & 12\\
    -3 & 15 & 6\\
    \end{pmatrix}$
\end{ejem}

Con esta definici\'on, es posible mostrar que: 

\begin{enumerate}
    \item $0A=0_{m\times n}$ para toda $A\in \M_{m\times n}(K)$.
    \item $\alpha(\beta A)=(\alpha \beta)A$ para toda $\alpha, \beta\in K$, para toda $A\in \M_{m\times n}(K)$.   
    \item $(\alpha+\beta )A=\alpha A +\beta A$ para toda $\alpha, \beta\in K$, para toda $A\in \M_{m\times n}(K)$.
    \item $\alpha(A + B) = \alpha A + \alpha B$ para toda $\alpha \in K$, para toda $A\in \M_{m\times n}(K)$.
    \item $1\cdot A=A$, $(-1)A=-A$ para toda $A\in \M_{m\times n}(K)$.
    \item $\tr(\alpha A)=\alpha \tr(A)$ para toda $A\in \M_{ n}(K)$.
    \item $(\alpha A)^t=\alpha A^t$ para toda $A\in \M_{m\times n}(K)$.
\end{enumerate}

\begin{rmk}
Toda matriz cuadrada $A$ se puede escribir como la suma de una matriz simétrica con otra
antisimétrica. Basta demostrar que para toda $A\in  \M_n $, ; la matriz $ A + A^t$ es simétrica y $A -A^t$
es antisimétrica, pues por la proposici\'on \ref{traspuesta} se tendría que
\[A = \frac{1}{2}(A + A^t)+\frac{1}{2}(A - A^t).\]
\end{rmk}

\begin{ejer}
Dado los siguientes comandos de \texttt{Octave}, muestre y argumente cada uno de los enunciados anteriores, proponga ejemplos para cada caso:
\vspace{0.5cm}
\begin{center}
\begin{tabular}{|>{\columncolor[gray]{0.8}} p{5cm}|p{10cm}|}
\hline
Comando & Explicación \\ \hline
\texttt{A=[1,2,3;4,5,6]}&  Define matriz de $2 \times 3$.\\ \hline
\texttt{A=zeros(m,n)}&  Define matriz nulas $m \times n$.\\ \hline
\texttt{B=ones(m,n)}&  Define matriz $m \times n$ llena de 1's.\\ \hline
\texttt{C=rand(m,n)}&  Define matriz $m \times n$ con números decimales al azar entre 0-1.\\ \hline
\texttt{D=randn(m,n)}&  Define matriz $m \times n$ con números de probabilidad normal $N(0,1)$. \\ \hline
\texttt{E=randp(p,m,n)} & Define matriz $m \times n$ con números enteros positivos de probabilidad poisson $P(\text{media}=p)$.\\ \hline
\texttt{I=eye(4,4)}&  Define matriz identidad $4 \times 4$.\\ \hline
\texttt{A’} & Transpuesta de $A$.\\ \hline
\texttt{A+B} & Suma de matrices\\ \hline
\texttt{3*A} & Multiplica el escalar 3 a cada elemento de $A$. \\ \hline
\texttt{-A} & Matriz inversa aditiva de $A$.\\ \hline
\texttt{T=trace(A)} & Devuelve el valor de la traza de $A$.\\ \hline
\texttt{V=diag(A)} & Devuelve en un vector vertical, los elementos de la
diagonal.\\ \hline
\texttt{Sum(V) o sum(A)} & Suma las coordenandas de un vector o suma las entradas columna a columna
de una matriz, entregando un vector.\\ \hline
\end{tabular}
\end{center}
\end{ejer}
\vspace{1.5cm}

\subsection{Ecuaciones matriciales}
Con las propiedades de las matrices podemos resolver ecuaciones algebraicas, es decir, buscamos una matriz tal que se cumpla una cierta igualdad.

\begin{ejer}
    Resuelva la ecuación y despeje la matriz $X$.
    
    \[5(X+A)-C=A-\frac{1}{3}(B-9X^t)^t,\]
    donde $A=\begin{pmatrix}
    2 & 1 \\
    -5 & 0 \\
    \end{pmatrix}$, $B=\begin{pmatrix}
    -3 & 6 \\
    1 & 2 \\
    \end{pmatrix}$ y $C^t=\begin{pmatrix}
    0 & 2 \\
    -1 & -2 \\
    \end{pmatrix}$.
\end{ejer}

\begin{ejer}
    Un fabricante produce tres modelos de zapatillas de descanso $A$, $B$ y $C$ en tres tamaños: para niños, damas y caballeros. La fabricación se realiza en dos plantas, una ubicada en
San Bernardo y la otra en Maipú. La producción semanal, en pares de zapatillas, en cada
planta se entrega a través de las matrices:



\begin{center}
\begin{tabular}{ >{\columncolor[gray]{0.8}}c c c c >{\columncolor[gray]{0.8}}c c c c}
San Bdo. &Niños& Damas& Varones& Maipú& Niños& Damas &Varones\\
A &20 &34 &30& A &16 &24& 26\\
B& 16 &20 &48 &B &10& 14& 32\\
C& 24 &28 &32& C& 15& 20 &28
\end{tabular}
\end{center}
\begin{enumerate}[(a)]
    \item Determine la matriz que contiene los datos relativos a la producción semanal total de cada modelo de zapatilla en ambas plantas.
    \item Si la producción en la planta de San Bernardo se incrementa en un 20\% y la de Maipú en un 40\%, escriba la matriz que representa la nueva producción semanal total de cada tipo de zapatilla.
    Desarrolle el problema con comandos de \texttt{Octave}.
\end{enumerate}
\end{ejer}

\subsection{Multiplicación de matrices}


\begin{defin}
Sean $a = \begin{pmatrix} a_1 \\ a_2 \\ \vdots \\ a_n \end{pmatrix}$ y $b = \begin{pmatrix} b_1 \\ b_2 \\ \vdots \\ b_n \end{pmatrix}$ dos vectores de igual dimensión. Entonces el producto escalar de $a$ y $b$ denotado por $a \cdot b$ está dado por 

$$a \cdot b = a_1b_1 + a_2b_2 + \cdots + a_nb_n$$
\end{defin}





\begin{ejem}
Sean $a =  \begin{pmatrix} -4 \\ -2 \\ 3 \end{pmatrix}$ y $b =  \begin{pmatrix} 3 \\ -2 \\ -5 \end{pmatrix}$. Entonces:

$$a \cdot b = \begin{pmatrix} -4 \\ -2 \\ 3 \end{pmatrix} \cdot \begin{pmatrix} 3 \\ -2 \\ -5 \end{pmatrix} = (-4)(3) + (-2)(-2) + (3)(-5) = -23$$
\end{ejem}

La multiplicación entre matrices $A \cdot B$ solo es posible cuando coincide el número de columnas
de $A$ con el número de filas de $B$, llamémosle a esta coincidencia $n$.
\begin{defin}
Sean $m, n, r$ n\'umeros naturales y sean $A=(a_{ij})\in \M_{m\times n}(K)$ y \newline ${B=(b_{ij})\in \M_{n\times r}(K)}$. La multiplicaci\'on de $A$ y $B$ es la matriz $AB=C=(c_{ij})$ de orden $m\times r$, donde 
\begin{align*}
   c_{ij}&=\sum^n_{k=1}a_{ik}b_{kj}.\\
 &= a_{i1}b_{1j} + a_{i2}b_{2j} + \dots +a_{in}b_{nj} 
\end{align*}
es decir,
\begin{center}
    $c_{ij} =$ (fila $i$ de $A$)$\cdot$(columna $j$ de $B$)
\end{center}
\end{defin}

\cmspecindice{Vector}
\cmspecindice{Matriz}

\begin{ejem}
Si $A=\begin{pmatrix}
    8 & -2 & -5\\
    3 & -1 & 1 \\
    \end{pmatrix}$ y $B=\begin{pmatrix}
    2 & -1 &1\\
    -2 & 4 &0\\
    1 &0 &0
    \end{pmatrix}$, entonces 
    \[AB=\begin{pmatrix}
    15 & -16 &8\\
    9 & -7 &3\\
    \end{pmatrix}\]
\end{ejem}


La multiplicación de matrices no es conmutativa, es decir, $AB \not= BA$. En el ejemplo anterior $BA$ no está definido porque la cantidad de columnas  de $B$ es distinta a la cantidad de filas de $A$.


\begin{ejem}
    $\begin{pmatrix}
   3 & 2 & \\
    1 & -1  \\
    \end{pmatrix}\cdot \begin{pmatrix}
    -4 & 3 \\
    2 & 1  \\
    \end{pmatrix}=\begin{pmatrix}
    -8 & 11 \\
    -6 & 2  \\
    \end{pmatrix}$, mientras que 
     $\begin{pmatrix}
    -4 & 3 \\
    2 & 1  \\
    \end{pmatrix}\cdot\begin{pmatrix}
   3 & 2 & \\
    1 & -1  \\
    \end{pmatrix} =\begin{pmatrix}
    -9 & -11 \\
    7 & 3  \\
    \end{pmatrix}$
\end{ejem}
En el caso particular en el cual $A\cdot B = B\cdot A$ en $\M_n$, diremos que $A$ y $B$ conmutan.

\begin{ejem}
    \item Un ejemplo de uso sencillo de la representación matricial es el de las matrices de inventario que abrevian una tabla bidimensional de doble entrada :
    
Si una multitienda tiene cuatro sucursales I, II, III y IV en las que se venden los mismos k productos $P_1, P_2, \dots, P_k$, desea resumir el stock ( en unidades de venta) que mantiene inmovilizado en las bodegas de cada local, podrá usar una matriz como sigue:
\begin{center}
\begin{tabular}{|c|c|c|c|c|c|}
\hline
Productos & $P_1$ & $P_2$ & $P_3$& \dots & $P_k$\\ \hline
I & $a_{11}$& $a_{12}$ &  $a_{13}$&\dots   & $a_{1k}$\\ \hline
II &     $a_{21}$& $a_{22}$ & $a_{23}$&\dots  & $a_{2k}$\\  \hline
III &     $a_{31}$& $a_{32}$  &$a_{33}$& \dots & $a_{3k}$\\  \hline
 IV&    $a_{41}$& $a_{42}$  &$a_{43}$ &\dots & $a_{4k}$ \\ \hline
 \end{tabular}
 \end{center}
 Si se define la matriz columna de precios de costos por unidad para cada uno de los $k$
productos $P_1, P_2, \dots, P_k$ como: $C = \left(\begin{array}{ccccc}
   c_1 & c_2 &  \dots& c_k 
\end{array}\right)^t$ y llamamos $S$ a la matriz de stock arriba
construida, entonces se podría obtener la matriz columna de capital inmovilizado por
sucursal:

\[M=S\cdot C=\begin{pmatrix}
a_{11}& a_{12} &  a_{13}&\dots   & a_{1k}\\
a_{21}& a_{22} & a_{23}&\dots  & a_{2k}\\
a_{31}& a_{32}  &a_{33}& \dots & a_{3k}\\
a_{41}& a_{42} &a_{43} &\dots & a_{4k}\\
\end{pmatrix}\begin{pmatrix}
   c_1 \\ c_2 \\  \vdots\\ c_k 
\end{pmatrix}=\begin{pmatrix} a_{11}c1+ a_{12}c_2 +  a_{13}c_3 + \dots   + a_{1k}c_k \\
\vdots\\
a_{41}c1+ a_{42}c_2 +  a_{43}c_3 + \dots   + a_{4k}c_k  \end{pmatrix}\]
que representa el costo total retenido en las bodegas de las sucursales I, II, III y IV.

Note que se definió una operación entre matrices para obtener $M$. Note además que
si semanalmente se registra una matriz $V$ de unidades vendidas por producto en cada
sucursal, al restar componente a componente la matriz $S$ con la matriz $V$, $S-V$ representaría
el remanente que queda en bodega al final de la semana. Si se le suma una matriz $A$ de
abastecimiento por producto y sucursal, se obtendrá una matriz $(S-V) + A$ con que se inicia
el nuevo stock para la semana siguiente.

Esto muestra que, si se definen operaciones apropiadas y de utilidad práctica, será
posible construir estructuras algebraicas de Matrices que resuelvan problemas de
contingencia como el mentado.
\end{ejem}

Ya habiendonos familiarizado con la operaci\'on producto, estudiaremos sus propiedades.

\begin{ejer}
    ¿Existe el elemento neutro para el producto de matrices? ¿Cu\'al?
\end{ejer}

Tal como en el caso de la suma, el producto de matrices tiene las siguientes propiedades:
\begin{thm}[Asociatividad]
Sean $m,n,p,q \in \N$. Sea $A$ una matriz de $n \times m$, $B$ de $m \times p$ y $C$ de $p \times q$. Entonces el producto es asociativo, es decir,
$$A(BC) = (AB)C.$$
\end{thm}
\begin{thm}[Distributividad]

Si todas las sumas y todos los productos siguientes están definidos (ejercicio: determine los ordenes de cada matriz, en ambos casos), entonces 
$$A(B+C) = AB+AC$$  $$(A+B)C = AC + BC$$
\end{thm}




Finalmente, tal como en el caso de la adici\'on de matrices, la matriz traspuesta y el producto, se relacionan de la siguiente manera:
\begin{thm}
Suponga que $A$ es una matriz de $n \times m$ y $B$ una matriz de $m \times p$.

\begin{enumerate}
\item
$(AB)^t = B^t A^t$.
\item $\tr(AB)=
\tr(BA)$ si $A, B\in \mathcal{M}_n(\mathbb{R})$.
\end{enumerate}
\end{thm}

En el producto de matrices no solo falla la conmutatividad, sino que adem\'as las potencias no se comportan como nuestra intuici\'on esperar\'ia. Para abordar estos casos, primero definiremos la potencia de una matriz:

\begin{defin}
Para una matriz cuadrada de orden $n$ definimos $A^2:=A\cdot A$, m\'as aun definimos $A^n:=A^{n-1}\cdot A$.
\end{defin}

\begin{rmk}
Note que si $A=(a_{ij})$, entonces $A^2\neq(a_{ij}^2).$
\end{rmk}
\begin{ejer}
    Sean $A, B$ matrices cuadradas de  orden $n$, ¿ser\'a siempre cierto que  $(A+B)^2=A^2+2AB+B^2$? Si la respuesta es negativa, ¿en qu\'e casos lo ser\'a?
\end{ejer}

\begin{defin}
Una matriz cuadradas $A$, no nulas tal que $A^2 = O$ se llaman \textbf{nilpotentes}. Diremos que la matriz es \textbf{nilpotente de orden $p$} si $ A^p= O$.
\end{defin}
Más aún, existen matrices $A$ y $B$, de orden $n$, no nulas, distintas tales que $AB = O$, es decir, el producto
de dos matrices puede ser la matriz cero y ninguna de ellas ser cero.
Por lo tanto: $AB = 0$ no es cierto que $A = 0 $ o $ B = 0$, a diferencia de los n\'umeros reales.

\begin{ejer}
    Muestre que la matriz $A=\begin{pmatrix}
    0 & -8 &0\\
    0 & 0  &0\\
    0& 5&0
    \end{pmatrix}$ es nilpotente.
\end{ejer}

\begin{defin}
\textbf{Idempotente de orden $p$} se llamará a la matriz que cumpla con $A^p = A$
\end{defin}

\begin{ejer}
    Muestre que las matrices $A=\begin{pmatrix}
    4 & -2 \\
    6 & -3  \\
    \end{pmatrix}$ y $B=\begin{pmatrix}
    2 & -3 &-5\\
    -1 & 4  &5\\
    1& -3&-4
    \end{pmatrix}$ son idempotentes.
\end{ejer}

\begin{defin}
\textbf{Matriz involutiva}, es aquella que $A^2=I_n$
\end{defin}

%\begin{ejer}
%    \begin{enumerate}
%        \item Sean $a, b$ y $c$ n\'umeros reales tales que $a^2+bc=a$, ¿qu\'e tipo de matriz es $C=\begin{pmatrix}
%    a & b \\
%    c & 1-a  \\
%    \end{pmatrix}$.
%    \item Encuentre una 
%    \end{enumerate}
%\end{ejer}


\subsection{Inversa de una matriz}

\begin{defin}
Diremos que una matriz $A\in\M_n$ es \textbf{invertible} si existe una matriz
$B \in \M_n$ tal que $A\cdot B = B\cdot A = I_n$ . Esta matriz $B$ se llamará \textbf{matriz inversa de $A$} y la denotaremos por $A^{-1}$. Luego usando cuantificadores:  
\begin{center}
  $A\in \M_n$ es invertible $\Longleftrightarrow$ $\exists \;B \in\M_n$ tal que $A\cdot B = B\cdot A = I_n$.  
\end{center}

\end{defin}

\begin{ejem}
Sea $A = \begin{pmatrix} 2 & 5 \\ 1 & 3 \end{pmatrix}$, entonces $A^{-1} = \begin{pmatrix} 3 & -5 \\ -1 & 2 \end{pmatrix}$.


\[A A^{-1} =  \begin{pmatrix} 2 & 5 \\ 1 & 3 \end{pmatrix}\begin{pmatrix} 3 & -5 \\ -1 & 2 \end{pmatrix} = \begin{pmatrix} 1 & 0 \\ 0 & 1 \end{pmatrix}\]


\[A^{-1}A = \begin{pmatrix} 3 & -5 \\ -1 & 2 \end{pmatrix}  \begin{pmatrix} 2 & 5 \\ 1 & 3 \end{pmatrix} =\begin{pmatrix} 1 & 0 \\ 0 & 1 \end{pmatrix} \]
\end{ejem}



%\begin{defin}
%Una matriz cuadrada que no es invertible se le denomina \textbf{singular} y una matriz invertible se denomina \textbf{no singular}.
%\end{defin}

\begin{thm}
\begin{enumerate}[(a)]
Sean $A$ y $B$ dos matrices invertibles de orden $n \times n$. Entonces:
    \item La  inversa de $A$ es única.
    
    \item $A^{-1}$ es invertible y $(A^{-1})^{-1} = A$.

    \item $AB$ es invertible y $(AB)^{-1} = B^{-1}A^{-1}$.

%    \item El sistema $Ax = b$ tiene una solución única $x = A^{-1}b$. Adem\'as, si $C$ es del mismo orden de $A$ y $AC=0$, entonces $C=0_n$. 

    \item Si $c$ un escalar distinto de cero, entonces $cA$ es una matriz invertible y $(cA)^{-1} = \frac{1}{c} A^{-1}$.

    \item $A^t$ es invertible y $(A^t)^{-1} = (A^{-1})^t$.

    \item $A^n$ es invertible para todos los enteros no negativos $n$ y $(A^n)^{-1} = (A^{-1})^n$.
\end{enumerate}
\end{thm}

\begin{rmk}
\begin{enumerate}
    \item Se puede ampliar la idea de matriz inversa para matrices no cuadradas, se
hablará aquí de inversas laterales: $A$ es invertible por la izquierda cuando exista $B \in
\M_{n x m}$ si $BA = I_n$ y $B$ será la Matriz inversa de $A$, por su izquierda. Diremos que $A$ es invertible por la derecha cuando exista $C \in
\M_{n x m}$ si $AC = I_m$ y $C$ será la Matriz inversa de $A$, por su derecha.
Mediante ejemplos concretos verifique que las inversas laterales puden no existir
o, de existir, no ser únicas.


\item En el caso de orden dos, es posible demostrar que una matriz del tipo
\[A=\begin{pmatrix}  a & b \\ c & d \end{pmatrix},\]
con $a, b,c,d \in \R$, es invertible si y solo si $ad-bc\neq 0$, en cuyo caso la inversa est\'a dada por
\[A^{-1}=\frac{1}{ad-bc}\begin{pmatrix}  d & -b \\ -c & a \end{pmatrix}.\]
\end{enumerate}
\end{rmk}

\begin{ejer}
\begin{enumerate}
\item
Sea $A = \begin{pmatrix}  2 & -3 \\ -4 & 5 \end{pmatrix}$. Calcule $A^{-1}$ si existe.
%\vspace*{5cm}

\item
Sea $A = \begin{pmatrix}  1 & 2 \\ -2 & -4 \end{pmatrix}$. Calcule $A^{-1}$ si existe.
%\vspace*{5cm}
\end{enumerate}
\end{ejer}
\vspace*{2cm}
\subsection{Resoluci\'on de ecuaciones matriciales usando matriz inversa}

Usando todas la propiedades y herramientas aprendidas, podemos resolver ecuaciones matriciales de manera m\'as eficiente.

\begin{ejem}
Sean $A = \begin{pmatrix}  1 & 1 & 0 \\ 0 & 2 & -1\\  0 & 0 & 1 \end{pmatrix}$ y $B = \begin{pmatrix}  3 & 0 & 1 \\ 0 & -1 & 1\\  1 & 1 & 0 \end{pmatrix}$. Despeje $X$ en la ecuaci\'on 
\[AX-3I=B.\]
Podemos observar, siempre que un escalar o constante sume o reste a una matriz deberá hacerlo
multiplicando con la identidad o con otra matriz, porque son objetos algebraicos de distinta naturaleza.
Primero notemos que la matriz $A$ es invertible y su inversa est\'a dada por 
\[A^{-1} = \begin{pmatrix}  1 & -1/2 & -1/2 \\ 0 & 1/2 & 1/2\\  0 & 0 & 1 \end{pmatrix},\]
(ejercicio: verifique). Luego
\begin{align*}
    AX-3I&=B\\
    AX&=B+3I /( A^{-1}\cdot) \text{ por la izquierda}\\
    (A^{-1}A)X&=A^{-1}(B+3I)\\
    IX&=A^{-1}(B+3I)\\
    X&=A^{-1}(B+3I)
\end{align*}
como 
\[B+3I=\begin{pmatrix}  3 & 0 & 1 \\ 0 & -1 & 1\\  1 & 1 & 0 \end{pmatrix}+\begin{pmatrix}  3 & 0 & 0 \\ 0 & 3 & 0\\  0 & 0 & 3 \end{pmatrix}=\begin{pmatrix}  6 & 0 & 1 \\ 0 & 2 & 1\\  1 & 1 & 3 \end{pmatrix},\]
obtenemos que 
\[X=A^{-1}(B+3I)=\begin{pmatrix}  1 & -1/2 & -1/2 \\ 0 & 1/2 & 1/2\\  0 & 0 & 1 \end{pmatrix}\begin{pmatrix}  6 & 0 & 1 \\ 0 & 2 & 1\\  1 & 1 & 3 \end{pmatrix}=\begin{pmatrix}  11/2 & -3/2 & -1 \\ 1/2 & 3/2 & 2\\  1 & 1 & 3 \end{pmatrix}\]
\end{ejem}

\begin{ejer}
    \begin{enumerate}
        \item Dado los siguientes comandos de \texttt{Octave}, muestre y argumente cada uno de los enunciados
anteriores, proponga ejemplos para cada caso:

\vspace{0.5cm}
\begin{center}
\begin{tabular}{|>{\columncolor[gray]{0.8}} p{5cm}|p{10cm}|}
\hline
Comando & Explicación \\ \hline
\texttt{[m,n]=size(A)}&  Asigna en $m$ y $n$ las dimensiones de $A$.\\ \hline
\texttt{A*B}&  Multiplicación de matrices.\\ \hline
\texttt{A*B}&  Realiza solo una multiplicación entre elementos $A$ con $B$, con la misma posición.\\ \hline
\texttt{A.+2}&  Suma 2 a cada elemento de $A$.\\ \hline
\texttt{inv(A)}&  Inversa de $A$. \\ \hline
\texttt{A/B} & Equivalente a: $A*$inv$(B)$.\\ \hline
\end{tabular}
\end{center}
    \item Desarrolle los pasos código a código de \texttt{Octave}, para resolver las ecuaciones:
    \begin{enumerate}
        \item Determine las matrices $A$ y $B$ que satisfacen el sistema:
        
        \begin{align*}
            2A+B&=\begin{pmatrix}  1 & 2 & 2 \\ -2 & 1 & 0 \end{pmatrix}\\
            A-3B&=\begin{pmatrix}  -4 & -3 & -2 \\ -1 & 0 & 1 \end{pmatrix}
        \end{align*}
        \item Despeje la matriz $X$ en la ecuación:
        \[AX^t+2B=(C-X)^t,\]
        donde $A=\begin{pmatrix}  1 & 0 \\ - 1 & 1 \end{pmatrix}$, $B=\begin{pmatrix}  0 & 2 \\ - 1 & -2 \end{pmatrix}$ y $C=\begin{pmatrix}  2 & 1 \\ 1 & 3 \end{pmatrix}$.
    \end{enumerate}

    \end{enumerate}
\end{ejer}
\newpage
\subsection{Operaciones elementales (escalonamientos) y matriz escalonada}

A partir de una matriz dada, el objetivo de esta secci\'on es obtener una matriz m\'as sencilla, pero que tenga las mismas propiedades que la matriz inicial.


Para una matriz $ A \in \M_{m\times n} (K)$ se definen, respecto de las filas de $A$, las siguientes \textbf{operaciones elementales}:
\begin{itemize}
    \item Permutar dos filas: Si se permuta la fila $i$ con la fila $j$ anotamos: $F_{i j}$ o $F_i\longleftrightarrow$.
    \item Multiplicar una fila por un número real: Si se multiplica la fila $i$ por el número $\alpha \neq 0$, esta operación
se anota: $F_i\longrightarrow\alpha  F_i$.
    \item Sumar dos filas: Si la fila $i$ se suma a la fila $j$, anotamos $F i + F j$, donde el resultado queda en la fila $j$.
    \item Finalmente, podemos combinar las dos últimas operaciones y obtener $F_j\longrightarrow \alpha  F_i + F_j$. El resultado queda en la fila $j$.
\end{itemize}

\begin{ejem}
\begin{align*}
     \begin{pmatrix}  2 & 1 & -3  \\ 5  & -4 & 1  \\ 1 & -1 & -4  \end{pmatrix} & F_3 \rightarrow F_1 = \begin{pmatrix}  1 & -1 & -4  \\ 5  & -4 & 1  \\ 2 & 1 & -3 \end{pmatrix} \\
     & F_3 \rightarrow F_2 =  \begin{pmatrix}  1 & -1 & -4  \\ 2 & 1 & -3  \\ 5  & -4 & 1  \end{pmatrix} \\
     &F_3 \rightarrow F_3 - 5F_1 = \begin{pmatrix} \ 1 & -1 & -4  \\ 2 & 1 & -3 \\  0 & 1 & 21  \end{pmatrix} \\
     &F_2 \rightarrow F_2 - 2F_1 = \begin{pmatrix}  1 & -1 & -4  \\ 0 & 3 & 5  \\  0 & 1 & 21  \end{pmatrix} \\
     &F_3 \rightarrow F_3 - \frac{1}{3} F_2 = \begin{pmatrix}  1 & -1 & -4  \\ 0 & 3 & 5  \\  0 & 0 & \frac{58}{3}  \end{pmatrix}
\end{align*}

\end{ejem}


\begin{defin}
Si $A, B \in \M_{m\times n}(K)$ y si $B$ se obtiene realizándole a la matriz $A$ un número finito de operaciones
elementales, entonces se dice que $A$ es equivalente a $B$ y se anota $A \sim B$.
\end{defin}

\begin{defin}
Se dice que $E$ es una matriz es escalonada cuando satisface las siguientes condiciones:
\begin{enumerate}
    \item Las filas que tengan todas sus componentes iguales a cero deben estar ubicadas debajo
de aquellas que tengan componentes no nulas.
\item La primera componente no nula de cada fila no nula es 1, vista de izquierda a derecha.
\item El número de ceros al comienzo de una fila aumenta a medida que se desciende en la
matriz.

\end{enumerate}

Si adem\'as, todas las componentes de la columna donde aparece un 1 distinguido son ceros, diremos que la matriz es escalonada reducida.
\end{defin}

\begin{ejem}
Las matrices $D=\begin{pmatrix}  1 & 0 & 0& 4\\ 0 & 1 & 0& 7\\
0 & 0 & 1& 2\\
0 & 0 & 0& 0\end{pmatrix}$ y $E=\begin{pmatrix}  1 & 0 & 5& 0\\ 0 & 1 & -2& 0\\
0 & 0 & 0& 1\end{pmatrix}$ son escalonadas.
\end{ejem}

\begin{ejer}
    Apliquemos operaciones elementales para obtener la escalonada de la matriz:
    \[\begin{pmatrix}  -2 & 0 & 1\\ 0 & 4 & -1\\1 & -2 & 0\end{pmatrix}\]
\end{ejer}
%\vspace*{10cm}

\vspace*{2cm}

\subsection{Rango de una matriz}

\begin{defin}
Si $A \in \M_{m\times n}(K)$, entonces $A \sim E$ , con $E$ matriz escalonada reducida por filas. Se define el
rango de $A$ como el número de filas no nulas de $E$. El rango de $A$ lo denotaremos por $r(A)$ o por $rank(A)$.
\end{defin}

\begin{rmk}
El rango de la matriz nula es cero $r(0_n) = 0$.\\
El rango de la matriz identidad es $r(I_n) = n$.
\end{rmk}

\begin{thm}
Sea $A\in  \M_n (K)$ entonces $A$ invertible si y solo si $r (A) = n$, es decir $A \sim I_n$
\end{thm}
\subsection{Procedimiento para encontrar la inversa de una matriz cuadrada $A$:}

\begin{enumerate}
\item
Se escribe la matriz aumentada $\begin{pmatrix}
A \bigm|& I
\end{pmatrix}$.

\item
Se utiliza la reducción por filas para poner la matriz $A$ a su forma escalonada reducida por filas.

\item
Se decide si $A$ es invertible:

\begin{enumerate}
\item
Si la forma escalonada reducida por filas de $A$ es la matriz identidad $I$, entonces $A^{-1}$ es la matriz que se tiene a la derecha de la barra vertical.

\item
Si la reducción de $A$ conduce a una fila de ceros a la izquierda de la barra vertical, entonces $A$ no es invertible pues su rango es menor a $n$.
\end{enumerate}
\end{enumerate} 

\begin{ejem}
\begin{align*}
     \begin{pmatrix}  2 & 1 & -3 &\bigm| 1&0 &0  \\ 5  & -4 & 1 &\bigm| 0&1 &0\\ 1 & -1 & -4  &\bigm| 0&0 &1\end{pmatrix} & F_3 \rightarrow F_1 = \begin{pmatrix}  1 & -1 & -4 &\bigm| \hspace*{3cm} \\ 5  & -4 & 1 &\bigm| \hspace*{3cm} \\ 2 & 1 & -3 &\bigm| \hspace*{3cm}\end{pmatrix} \\
     & F_3 \rightarrow F_2 =  \begin{pmatrix}  1 & -1 & -4 &\bigm| \hspace*{3cm} \\ 2 & 1 & -3 &\bigm| \hspace*{3cm} \\ 5  & -4 & 1 &\bigm| \hspace*{3cm} \end{pmatrix} \\
     &F_3 \rightarrow F_3 - 5F_1 = \begin{pmatrix} \ 1 & -1 & -4 &\bigm| \hspace*{3cm}  \\ 2 & 1 & -3 &\bigm| \hspace*{3cm}\\  0 & 1 & 21 &\bigm| \hspace*{3cm}  \end{pmatrix} \\
     &F_2 \rightarrow F_2 - 2F_1 = \begin{pmatrix}  1 & -1 & -4 &\bigm| \hspace*{3cm} \\ 0 & 3 & 5 &\bigm| \hspace*{3cm} \\  0 & 1 & 21 &\bigm| \hspace*{3cm} \end{pmatrix} \\
     &F_3 \rightarrow F_3 - \frac{1}{3} F_2 = \begin{pmatrix}  1 & -1 & -4 &\bigm| \hspace*{3cm} \\ 0 & 3 & 5 &\bigm| \hspace*{3cm} \\  0 & 0 & 58/3 &\bigm| \hspace*{3cm} \end{pmatrix}\\
     &F_2 \rightarrow \frac{1}{3}F_2 = \begin{pmatrix}  1 & -1 & -4 &\bigm| \hspace*{3cm} \\ 0 & 1 & 5/3 &\bigm| \hspace*{3cm} \\  0 & 0 & 58/3 &\bigm| \hspace*{3cm} \end{pmatrix}\\
    &F_3 \rightarrow \frac{3}{58}F_3 = \begin{pmatrix}  1 & -1 & -4 &\bigm| \hspace*{3cm} \\ 0 & 1 & 5/3 &\bigm| \hspace*{3cm} \\  0 & 0 & 1 &\bigm| \hspace*{3cm} \end{pmatrix}\\
        &F_2 \rightarrow F_2 -\frac{5}{3}F_3 = \begin{pmatrix}  1 & -1 & -4 &\bigm| \hspace*{3cm} \\ 0 & 1 & 0 &\bigm| \hspace*{3cm} \\  0 & 0 & 1 &\bigm| \hspace*{3cm} \end{pmatrix}\\
                &F_1 \rightarrow F_1 +F_2+ 4F_3  = \begin{pmatrix}  1 & 0 & 0 &\bigm| \hspace*{3cm} \\ 0 & 1 & 0 &\bigm| \hspace*{3cm} \\  0 & 0 & 1 &\bigm| \hspace*{3cm} \end{pmatrix}
\end{align*}
\begin{ejer}
    \begin{enumerate}
        \item Determine todos los valores de a para que $r(M)=3$, con
            \[\begin{pmatrix}  1 & 2 & a\\ 2 &3 & 1\\3 & 4 & a^2\end{pmatrix}\]
        \item Realice operaciones elementales entre filas para obtener la escalonada e indique el rango de la matriz
dependiendo de los valores reales de $k$:
            \[\begin{pmatrix}  1 & -1 & 0 & 1\\ 1 &k+1 & 1 & 0\\1 & -4 & 3& -2\end{pmatrix}\]
        \item Dado los siguientes comandos de Octave, muestre y argumente cada uno de los enunciados
anteriores, proponga ejemplos para cada caso:
\vspace{0.5cm}
\begin{center}
\begin{tabular}{|>{\columncolor[gray]{0.8}} p{5cm}|p{10cm}|}
\hline
Comando & Explicación \\ \hline
\texttt{rank(A)}&  Calcula el rango de $A$.\\ \hline
\texttt{rref(A)}&  Determina la matriz escalonada de $A$, realiza las
operaciones elementales entre fila.\\ \hline
%\texttt{rref([A;eye(3)])}&  Si $A$ es invertible, calcula así la inversa conoperaciones elementales.\\
\hline

\end{tabular}
\end{center}
Enlace recomendado para operaciones con literales: \url{https://es.symbolab.com/solver/matrix-calculator}.
    \end{enumerate}
\end{ejer}
\end{ejem}

